%! Author = matteomagnini
%! Date = 06/05/26

\documentclass[11pt]{article}

\usepackage{amsmath}
\usepackage{amssymb}
\usepackage{hyperref}
\usepackage{tikz}

\title{
    Intelligent System 2\\
    Tutorial week 12: Structured Argumentation
}

\author{
    Matteo Magnini\\
    \href{mailto:matteo.magnini@uni.lu}{matteo.magnini@uni.lu}
}

\date{
    13 May, 2026
}

\begin{document}

\maketitle

\section{Exercises}

\subsection{Exercise 1: Formalizing a Knowledge Base}

Consider the following scenario:

\begin{itemize}
    \item The road looks wet.
    \item If the road looks wet, then typically the road is wet.
    \item If the road is wet, then typically driving is dangerous.
    \item There is a mirage effect.
    \item If there is a mirage effect, then typically the road is not wet.
\end{itemize}

Use the atoms:

\[
roadLooksWet,
roadWet,
dangerous,
mirage
\]

\begin{enumerate}
    \item Identify:
    \begin{itemize}
        \item facts
        \item strict rules
        \item defeasible rules
    \end{itemize}

    \item Write the knowledge base:
    \[
    K = (R_S,R_D)
    \]
\end{enumerate}

\subsection{Exercise 2: Building Arguments}

Using the knowledge base from Exercise 1:

\begin{enumerate}
    \item Construct the arguments:
        \begin{itemize}

            \item textit{The road looks wet.}

            \item \textit{Given that the road looks wet, it is typically the case that the road is wet.}

            \item \textit{Given that there is a mirage effect, it is typically the case that the road is not wet.}

            \item \textit{Given that the road is wet, it is typically the case that driving is dangerous.}
        \end{itemize}

    \item For each argument specify:
    \begin{itemize}
        \item conclusion
        \item subarguments
        \item top rule
    \end{itemize}
\end{enumerate}

\subsection{Exercise 3: Computing Attacks}

Consider the arguments from Exercise 2.

\begin{enumerate}
    \item Which arguments rebut each other?

    \item Which attacks propagate to superarguments?

    \item Draw the abstract argumentation framework.

    \item Is:
    \[
    \{A_3,A_4\}
    \]
    conflict-free?
\end{enumerate}

\subsection{Exercise 4: Weakest Link and Last Link}

Suppose:

\[
str(r_1)=1
\]

\[
str(r_2)=3
\]

where:

\[
r_1: roadLooksWet \Rightarrow roadWet
\]

\[
r_2: roadWet \Rightarrow dangerous
\]

Consider:

\[
A=[roadLooksWet] \Rightarrow roadWet
\]

\[
B=[A] \Rightarrow dangerous
\]

\begin{enumerate}
    \item Compute the strength of:
    \[
    A
    \]
    and
    \[
    B
    \]
    using:
    \begin{itemize}
        \item Weakest Link
        \item Last Link
    \end{itemize}

    \item Which principle gives more importance to the final rule?
\end{enumerate}

\end{document}
